% !TEX root = ../../main-jos.tex
\section{实验与分析}

\subsection{环境与方法}
实验在 WSL2（8\,GB 内存，4\,vCPU）Docker 原型集群上进行。集群含 1 个 OpenResty 网关、8 个 httpbin 微服务及对应 Sidecar Agent、1 个日志中心及 Loki/Redis 等基础设施。应用日志采用 Apache Combined 风格字段，URL 与网关路由一致。

对比实验将定向采集与全量采集（\texttt{COLLECTION\_MODE=full}）在同一硬件环境下对照：负载为 180\,s、并发 50 的混合 HTTP 请求（含正常/500 错误/慢请求），每模式独立重复 3 次。Loki 入库增量取自 Center 统计接口；Agent CPU/内存于每轮压测结束后对全部 9 个 Agent 容器执行 \texttt{docker stats} 单次快照并取算术均值。每轮网关流量日志 Redis 增量约 $1.97\times10^4$ 条，关注清单稳定生成 16 条 URL 模式（含 \texttt{/status/500}、\texttt{/delay/1} 等异常/慢请求泛化模式），保障高价值请求可被定向规则覆盖。原始数据见 phase3 结果 JSON。

\subsection{研究问题与指标}
为使实验与贡献形成可审计对应关系，本文围绕 4 个研究问题组织实验，如表~\ref{tab:rq_mapping} 所示。RQ1 对应采集前端减量能力，RQ2 对应 Sidecar 资源开销，RQ3 对应三层策略贯通能力，RQ4 对应缓存与退避机制的可靠性边界。

\begin{table*}[htbp]
  \centering
  \caption{研究问题、指标与证据映射}
  \label{tab:rq_mapping}
  \small
  \resizebox{\textwidth}{!}{%
  \begin{tabular}{p{1.2cm}p{5.2cm}p{4.2cm}p{4.0cm}}
    \toprule
    编号 & 研究问题 & 指标 & 证据位置 \\
    \midrule
    RQ1 & 定向采集能否显著降低 Loki 入库日志量？ & Loki 入库条数、带宽估计、降低比例 & 表~\ref{tab:compare}、图~\ref{fig:compare} \\
    RQ2 & 定向采集是否控制 Sidecar 资源开销？ & Agent CPU、内存均值与标准差 & 表~\ref{tab:compare}、图~\ref{fig:compare} \\
    RQ3 & 网关预筛选、节点定向采集与中心二次过滤是否形成端到端闭环？ & Redis 网关日志增量、关注清单模式数、链路阶段打通情况 & §6.3 与原始 phase3 JSON \\
    RQ4 & 固定缓存块与压力感知退避是否具有必要性与有界性？ & 策略仿真消融、微基准耗时、退避单元测试与等待上界 & 图~\ref{fig:ablation}、图~\ref{fig:attention}、式~\ref{eq:backoff} \\
    \bottomrule
  \end{tabular}%
  }
\end{table*}

\subsection{实验结果}
表~\ref{tab:compare} 与图~\ref{fig:compare} 给出 RQ1 与 RQ2 的三期实测对比。相较同架构全量模式，定向采集在 Loki 入库日志量上降低 98.4\%（72$\pm$0 条 vs.\ 4388$\pm$1 条），在 Agent CPU 上相对降低 37.5\%（绝对值分别为 0.05$\pm$0.01\% 与 0.08$\pm$0.02\%，标准差存在重叠，宜结合效应量解读）。需指出：全量基线除关闭规则过滤外，应用日志源发射频率（400\,ms）亦高于定向模式（2\,s），表~\ref{tab:compare} 降幅为\textbf{策略过滤与源配置差异的联合效果}。

\textbf{发射频率敏感性。}源端理论产生速率比约为 $2000/400{=}5{:}1$。若入库降幅仅由发射频率导致，预期入库比应接近 5:1；实测为 $4388/72{\approx}61{:}1$，表明\textbf{定向过滤策略}贡献了绝大部分减量。因此 98.4\% 降幅以策略过滤为主、发射频率差异为辅；同频率对照实验列入 §6.6 后续工作。即便考虑该因素，定向模式仍将入库量控制在百条量级。

\textbf{成本粗算。}按表~\ref{tab:compare} 带宽均值，单场 180\,s 压测南北向日志流量约减少 $(12.19{-}0.20){\times}180{\approx}2.2$\,MB；在类似负载持续运行下，定向模式可显著降低上行带宽与 Loki 索引写入压力（具体节约量取决于 QPS 与日志字段大小）。图~\ref{fig:ablation} 为 RQ4 的策略仿真消融（非集群实测）。

\begin{table}[htbp]
  \centering
  \caption{定向采集与全量采集资源对比（八节点，180\,s，$n{=}3$ 实测）}
  \label{tab:compare}
  \begin{tabular}{lccc}
    \toprule
    指标 & 定向采集 & 全量采集 & 降低比例 \\
    \midrule
    Loki 入库 (条) & 72$\pm$0 & 4388$\pm$1 & 98.4\% \\
    Agent CPU (\%) & 0.05$\pm$0.01 & 0.08$\pm$0.02 & 37.5\% \\
    Agent 内存 (MB) & 95.6$\pm$4.2 & 97.5$\pm$1.7 & 2.0\% \\
    带宽 (KB/s) & 0.20 & 12.19 & 98.4\% \\
    \bottomrule
  \end{tabular}
\end{table}

\begin{figure}[htbp]
  \centering
  \includegraphics[width=0.85\textwidth]{fig3_comparison_bars.pdf}
  \caption{定向采集与全量采集资源消耗对比（八节点，180\,s，$n{=}3$ 实测）}
  \label{fig:compare}
\end{figure}

\begin{figure}[htbp]
  \centering
  \includegraphics[width=0.95\textwidth]{fig4_ablation_bars.pdf}
  \caption{消融实验结果（\textbf{策略仿真}，非集群实测）}
  \label{fig:ablation}
\end{figure}

针对 RQ3，端到端链路已实测打通：Nginx 共享内存 $\rightarrow$ \allowbreak Agent gRPC $\rightarrow$ \allowbreak Redis $\rightarrow$ \allowbreak 关注清单生成 $\rightarrow$ \allowbreak Agent 规则拉取（复现脚本见 \url{experiments/scripts/phase3_collect.py}，原始结果见 \url{experiments/results/phase3/phase3_latest.json}）。针对 RQ4，关注清单生成在 5000 条日志上耗时约 15\,ms（图~\ref{fig:attention}）；指数退避模块通过 \url{agent/pkg/retry/backoff_test.go} 全部单元测试用例。

\begin{figure}[htbp]
  \centering
  \includegraphics[width=0.85\textwidth]{fig6_attention_bench.pdf}
  \caption{关注清单生成算法微基准}
  \label{fig:attention}
\end{figure}

\subsection{可复现性与效度威胁}
实验脚本、Docker Compose 配置与原始 JSON 结果保存在 \texttt{experiments/} 及 phase3 结果目录，图~\ref{fig:compare}、表~\ref{tab:compare} 均由该 JSON 派生。当前结果仍存在以下效度威胁。

\textbf{规模效度}：八/十六节点为 WSL/Docker 实测，三十二/六十四节点为基于实测曲线的外推；云环境长稳与带宽压力仍需验证。

\textbf{基线效度}：已给出同架构全量基线与 Promtail/OTel/eBPF 同负载估算（图~\ref{fig:baseline}），完整端到端部署对照仍待补充。

\textbf{负载效度}：当前负载由 httpbin 合成，覆盖正常、错误与慢请求，但距离 DeathStarBench、Alibaba 生产迹等真实微服务流量仍有差距。

\textbf{统计效度}：三期对比为 180\,s、$n{=}3$；四期规模矩阵为 60\,s、$n{=}3$。逐探针命中率受 Agent 批处理影响，不能单独代表稳态漏报率。图~\ref{fig:ablation} 为策略仿真消融，不能与表~\ref{tab:compare} 的集群实测直接比较。

\subsection{规模扩展与质量指标}
在四期实验中，本文在八节点与十六节点集群上补跑定向模式规模矩阵（60\,s、并发 50、每规模重复 3 次），并对 32/64 节点给出基于实测曲线的资源外推。图~\ref{fig:scale} 显示：定向模式下 Loki 入库量随节点数缓慢增长（八节点 27$\pm$0 条、十六节点 47.7$\pm$5.8 条/轮），而 Agent 平均 CPU 维持在 0.02\%--0.03\% 量级；Center CPU 在十六节点下为 0.28$\pm$0.41\%，未出现随规模线性失控。32/64 节点柱形为外推值（灰色），供云环境长稳验证参考。

\begin{figure}[htbp]
  \centering
  \includegraphics[width=0.95\textwidth]{fig7_scale_bars.pdf}
  \caption{多规模扩展实验（绿：实测 8/16 节点；灰：32/64 外推）}
  \label{fig:scale}
\end{figure}

\textbf{高价值日志漏报与端到端延迟。}稳态混合负载下，关注清单稳定覆盖 \texttt{/status/500}、\texttt{/delay/1} 等泛化模式（八节点 16 条、十六节点 32 条），\textbf{规则级}漏报率为 0（合成负载与清单配置下）。生产迹风格混合负载（60\% 正常 / 20\% 5xx / 12\% 慢请求 / 8\% 其他，90\,s、十六节点）共发送 3.6$\times10^4$ 次请求，定向入库 77 条，验证减量能力。逐探针实验的即时命中率受 Agent 2\,s 批处理影响，不等同于规则级漏报，亦不能单独代表生产稳态漏报率。端到端延迟（HTTP $\rightarrow$ Center \texttt{received}）实测：P50 0.83\,s、P95 12\,s（$n{=}30$）；P95 较高与批处理窗口及合成慢请求有关，面向实时告警需结合业务 SLA 评估可接受性。中心故障恢复实验（\texttt{docker stop/start}）恢复时间约 1.1\,s。

\subsection{工业基线对照}
图~\ref{fig:baseline} 在同负载口径下给出本文定向采集与三类工业基线的对照。\textbf{仅 P3 定向采集为八节点三期集群实测}（72 条、CPU 0.05\%）；Promtail+静态过滤、OpenTelemetry 尾部采样和 eBPF 探针方案基于同架构全量基线与公开文献开销模型\textbf{估算}，\textbf{尚未在本文集群中完整部署}，图中柱形应解读为趋势参考而非实测结论。估算结果表明：定向采集在保持最低入库量的同时，CPU 开销与 Promtail 静态过滤同量级，显著低于尾采样 Collector 的推理与追踪附加成本；Promtail 同集群实测列入后续工作。

\begin{figure}[htbp]
  \centering
  \includegraphics[width=0.85\textwidth]{fig8_baseline_bars.pdf}
  \caption{工业基线对照（P3 实测，其余为同负载估算）}
  \label{fig:baseline}
\end{figure}

\subsection{剩余局限与后续工作}
四期实验已补充 8/16 节点规模矩阵、规则级漏报验证、端到端延迟与工业基线估算，但 32/64 节点仍为资源外推，Promtail/OpenTelemetry/eBPF 尚需端到端部署实测，同频率发射对照与集群实测消融尚待补充。后续将在云环境开展 32/64 节点长稳压测，引入 DeathStarBench 等真实负载，完成 Promtail 同集群实测、同频率对照实验，并将集群实测消融纳入投稿前证据包。
